\documentclass[12pt]{letter}
\usepackage[margin=1in]{geometry}
\usepackage{times}
\usepackage{hyperref}
\hypersetup{colorlinks=true, urlcolor=blue}

\signature{Sayuj Krishnan S., MS, MCh\\Department of Neurosurgery\\Yashoda Hospitals, Hyderabad 500003, India\\hellodr@drsayuj.info}

\address{Sayuj Krishnan S.\\Department of Neurosurgery\\Yashoda Hospitals\\Hyderabad 500003, India}

\begin{document}

\begin{letter}{The Editor\\Nature}

\opening{Dear Editor,}

We submit our manuscript entitled \textbf{``Metabolic buckling links growth scaling to scoliotic curve diversity''} for consideration as an Article in \textit{Nature}.

Adolescent idiopathic scoliosis (AIS) affects 3\% of children worldwide, yet after more than a century of investigation, its cause remains unknown. The defining puzzle is temporal: 87\% of scoliotic curves emerge within six months of peak height velocity --- a window spanning less than 2\% of skeletal development. No existing theory explains this precision, nor why the deformity is invariably three-dimensional, manifests in six distinct Lenke curve patterns, or disproportionately affects girls 10:1.

Our work resolves all four puzzles from a single physical principle. We demonstrate that the metabolic power required to maintain spinal alignment scales as $L^4$ (spinal length to the fourth power), while the metabolic supply during rapid growth scales as $L^2$. This creates a transient energy deficit window during the adolescent growth spurt that triggers sequential failure of the most metabolically expensive mechanosensory proteins --- a cascade we verify using structural analysis of 47 proteins via AlphaFold. The framework makes quantitative, falsifiable predictions for each Lenke curve type, the female prevalence, and the characteristic three-dimensionality of the deformity.

This work is significant for several reasons:
\begin{enumerate}
\item It provides the first \textit{unified} pathogenetic theory for AIS, subsuming rather than competing with established hypotheses (melatonin, genetics, circadian desynchrony).
\item It introduces a novel metric (Conformational Ordering Index) that links protein structural biology to disease susceptibility through thermodynamic principles.
\item It identifies metabolic rescue as a rational therapeutic strategy --- potentially transforming AIS management from reactive mechanical correction to proactive biochemical prevention.
\item The framework connects biomechanics, statistical physics, and developmental biology in a manner that will be of broad interest to Nature's interdisciplinary readership.
\end{enumerate}

This manuscript has not been submitted elsewhere and all authors have approved the submission. We declare no competing interests. Simulation code and AlphaFold analysis data are publicly available at \url{https://github.com/spinalmodes}.

We suggest the following potential reviewers with relevant expertise:

\begin{itemize}
\item \textbf{Prof.\ Stuart Weinstein} (University of Iowa) --- clinical AIS epidemiology
\item \textbf{Prof.\ Alain Moreau} (Universit\'{e} de Montr\'{e}al) --- AIS molecular biology
\item \textbf{Prof.\ M.\ Gail Faulkner} (University of Alberta) --- spinal biomechanics
\item \textbf{Prof.\ Geoffrey West} (Santa Fe Institute) --- allometric scaling in biology
\end{itemize}

We believe our findings will be of great interest to Nature's broad readership, bridging fundamental physics and one of the most common unsolved problems in paediatric orthopaedics.

\closing{Yours sincerely,}

\end{letter}
\end{document}
